% problem_id: S047-aif3605-N1
% source_id: S047 (Magnetization in the zig-zag layered Ising model and orthogonal polynomials)
% source_file: aif3605.pdf
% statement_locator: printed p. 2309 (printed; = PDF p. 36 of aif3605.pdf)
% proof_locator: printed p. ['2309 (printed; = PDF p. 36) — the proof is a single paragraph on the same page as the statement', "enabling context: 2307–2309 (printed, = PDF 34–36); the authors' own reading Remark 4.3: printed 2310 (= PDF 37); co-cited Lemma 4.1: printed 2305–2307 (= PDF 32–34)"]
% representation: PDF; transcribed from rendered page images (never from the text layer)
% collected_by: carrier rule (the headline theorem was an assembly; this is the step it relies on)
% review_status: PENDING independent review (single delegated judgement)
% status: complete (statement + author proof verbatim + reason + locators)
%
%% =====================================================================
%% human-proof-corpus / deepseek-takeover / batch b17 -- candidate N1
%% =====================================================================
%% Paper : D. Chelkak, C. Hongler, R. Mahfouf, "Magnetization in the
%%         zig-zag layered Ising model and orthogonal polynomials",
%%         Ann. Inst. Fourier (Grenoble) 74, no. 6 (2024), 2275-2330.
%%         doi:10.5802/aif.3605
%% Source: /disks/sata1/yupeng/human-proof-corpus/source-cache/
%%         registry-sources/numdam/aif3605.pdf   (57 PDF pages)
%%         page map: PDF page = printed page - 2273 (checked on pp. 32-38)
%% ---------------------------------------------------------------------
%% Original target (the "signboard" handed to me): Theorem 1.1,
%%   Introduction, printed p. 2279 (= PDF p. 6).
%%
%% COLLECTED THEOREM (the bearer): PROPOSITION 4.2, printed p. 2309
%%   (= PDF p. 36) -- in the authors' words "the key proposition on the
%%   construction of solutions to (4.2), (4.3) in the upper quadrant".
%%
%% WHY.  The printed proof of Theorem 1.1 is an assembly.  Its single
%%   load-bearing sentence is
%%      "Therefore, Lemma 4.1 and Proposition 4.2 imply that
%%       D_even H_0^o = D_even U_even H_0 = S_even H_0 = J^{1/2}[...]"
%%   (printed p. 2310); everything that follows is determinant linear
%%   algebra.  So the signboard Theorem 1.1 itself is H2 (a two-name
%%   reduction).  The step of the deferral that carries a SELF-MADE
%%   CONSTRUCTION is Proposition 4.2: it introduces the operator D
%%   (the authors: "an operator D, which plays the key role in the rest
%%   of the analysis"), the identities (4.9)
%%      CD = D C^o,   D D* = I - C^2,   D* D = I - (C^o)^2,
%%   and builds by polar decomposition the (partial) isometry
%%   U = (D*D)^{-1/2} D* -- the very operator U_even whose truncated
%%   determinant equals the magnetization M_m in (1.5).  This is the
%%   object that actually bears formula (1.5).
%%
%%   The deferral in the proof of Theorem 1.1 is a TWO-NAME deferral.
%%   The co-cited companion, Lemma 4.1 (uniqueness of the half-plane
%%   fermionic observable, printed pp. 2305-2307), is the rigidity input
%%   that lets the constructed pair be identified with the actual
%%   observable; it is the "routine of the method" in this
%%   (Chelkak-Smirnov) circle of ideas, and its full proof is
%%   nevertheless transcribed here in Appendix B so that the unit can be
%%   re-judged if the collection line prefers the probabilistic
%%   uniqueness lemma as the bearer.
%% ---------------------------------------------------------------------
%% Rendering: pdftoppm -f <P> -l <P> -r 150 -png aif3605.pdf <out>
%%   (a 300 dpi second pass was used for two ambiguous glyphs on
%%   pp. 2307, 2308; the wording is unaffected).
%% Transcription: verbatim LaTeX, nothing rewritten, nothing completed.
%%   Page breaks are marked "%% --- page NNNN ---" with the PRINTED page
%%   number.  Source oddities are copied as printed and flagged inline
%%   with "[as printed: ...]".
%% =====================================================================
\documentclass[11pt]{article}
\usepackage[margin=2.3cm]{geometry}
\usepackage{amsmath,amssymb,amsfonts}
\usepackage{enumitem}
\setlist{topsep=2pt,itemsep=1pt,parsep=0pt}
\newcommand{\eql}[1]{\text{(#1)}\quad}
\newcommand{\Dev}{D_{\mathrm{even}}}
\newcommand{\Uev}{U_{\mathrm{even}}}
\newcommand{\Sev}{S_{\mathrm{even}}}
\newcommand{\bHd}{\mathbb{H}^{\diamond}}
\newcommand{\Htil}{\widetilde{H}}
\newcommand{\Htilo}{\widetilde{H}^{\circ}}
\newcommand{\bvv}{\mathbf{v}}
\allowdisplaybreaks
\title{human-proof-corpus b17 / candidate N1\\[2pt]
\large transcription: Proposition 4.2 of aif3605 (the bearer of Theorem 1.1)}
\date{}
\begin{document}
\maketitle
\vspace{-3em}

\noindent\rule{\textwidth}{0.4pt}
\noindent\textbf{Reading key.} Block 0 = the original target (the signboard).
Block 1 = the prerequisites the collected unit quotes. Block 2 = the
enabling construction that the collected unit lives in. \textbf{Block 3 = the
collected unit (Proposition 4.2 and its proof).} Appendix A = the
assembly evidence (proof of Theorem 1.1). Appendix B = the co-cited
companion (Lemma 4.1 and its proof).
\noindent\rule{\textwidth}{0.4pt}

%% =====================================================================
\section*{[Block 0] Original target (signboard): Theorem 1.1 --- printed p.~2279}
%% =====================================================================

\noindent\emph{(Context sentence from printed p.~2278; verbatim.)}

The (half-)infinite volume limit of the Ising model on $\bHd$ is defined as a
limit of probability measures on an increasing sequence of finite domains
exhausting $\bHd$, with ``$+$'' boundary conditions at the right-most column
$\mathrm{C}_0$ and at infinity. All interaction parameters between the
columns $\mathrm{C}_{p-1}$ and $\mathrm{C}_p$ are assumed to be the same and
equal to $x_p = \exp[-2\beta J_p] = \tan\frac12\theta_p$, where
$\theta_p \in (0,\frac12\pi)$ can be viewed as a convenient parametrization of
$\beta J_p$, see Section 2.1 for more details. Let
\[
\eql{1.1} M_m = M_m(\theta_1,\theta_2,\dots) := \mathbb{E}^{+}_{\bHd}\bigl[\sigma_{(-2m-\frac12,0)}\bigr]
\]

%% --- page 2279 ---
be the magnetization in the $(2m)$-th column (the analysis for odd columns
can be done similarly). Denote
\[
\eql{1.2} \Dev := i
\begin{bmatrix}
\cos\theta_1\cos\theta_2 & 0 & 0 & \cdots\\
-\sin\theta_2\sin\theta_3 & \cos\theta_3\cos\theta_4 & 0 & \cdots\\
0 & -\sin\theta_4\sin\theta_5 & \cos\theta_5\cos\theta_6 & \cdots\\
\cdots & \cdots & \cdots & \cdots
\end{bmatrix}
\]
and let
\[
\eql{1.3} D_{\mathrm{even}}^{*} = \Uev \Sev,\qquad
\Sev = (\Dev D_{\mathrm{even}}^{*})^{1/2}
\]
be the polar decomposition of the operator $D_{\mathrm{even}}^{*}$, see also
Remark 4.3 for another interpretation of the (partial) isometry $\Uev$.
Further, denote $J := \Dev D_{\mathrm{even}}^{*}$. A straightforward
computation shows that
\[
\eql{1.4} J =
\begin{bmatrix}
b_1 & -a_1 & 0 & \cdots\\
-a_1 & b_2 & -a_2 & \cdots\\
0 & -a_2 & b_3 & \cdots\\
\cdots & \cdots & \cdots & \cdots
\end{bmatrix},
\]
where
\[
a_k = \cos\theta_{2k-1}\cos\theta_{2k}\sin\theta_{2k}\sin\theta_{2k+1},
\]
\[
b_k = \cos^2\theta_{2k-1}\cos^2\theta_{2k} + \sin^2\theta_{2k-2}\sin^2\theta_{2k-1},
\]
$\theta_0 := 0$ and $b_1 = \cos^2\theta_1\cos^2\theta_2$.

Let $\nu_J$ be the spectral measure of $J$ associated with the first basis
vector. It is easy to see that $0 \leqslant J \leqslant 1$ and thus
$\operatorname{supp}\nu_J \in [0,1]$. Given a measure $\mu$ on $[0,1]$, let
$H_m[\mu] := \det[\int_0^1\lambda^{p+q}\mu(d\lambda)]_{p,q=0}^{m-1}$ be the
$m$-th Hankel determinant composed from the moments of this measure. Denote
by $P_m$ the orthogonal projector on the space of first $m$ coordinates of
$\ell^2$.

\medskip
\noindent\textsc{Theorem 1.1.} --- \emph{For all $\theta_1,\theta_2,\dots \in
(0,\frac{\pi}{2})$ and $m \geqslant 1$, we have}
\begin{align*}
\eql{1.5} M_m &= |\det P_m \Uev P_m|\\
&= \frac{\det P_m J^{1/2}P_m}{\prod_{k=1}^{2m}\cos\theta_k}
= \frac{H_m[\lambda^{1/2}\nu_J]}{(H_m[\nu_J]\cdot H_m[\lambda\nu_J])^{1/2}},
\end{align*}
%% [the trailing comma of (1.5) is as printed.]
\emph{where $\Uev$ is the (partial) isometry factor in the polar
decomposition (1.3), the Jacobi matrix $J = \Dev D_{\mathrm{even}}^{*}$ is
given by (1.4), and $\nu_J$ is the spectral measure of $J$.}

%% =====================================================================
\section*{[Block 1] Prerequisites quoted by the collected unit --- printed p.~2305}
%% =====================================================================

\noindent\emph{(The sentence begins on printed p.~2304: ``The discrete
Cauchy--Riemann identities (2.10) can be written''; the page below starts
with its last word.)}

%% --- page 2305 ---
as
\[
\eql{4.2} H(-k-1,s\pm1)\sin\theta_{k+1} - H(-k,s)\cos\theta_k
= ,\pm i\cdot\bigl[H^{\circ}(-k,s\pm1)\sin\theta_k - H^{\circ}(-k-1,s)\cos\theta_{k+1}\bigr]
\]
%% [as printed: the source prints "= , ±i ·" (comma before the sign); copied verbatim.]
if $k \geqslant 1$ and $k+s \notin 2\mathbb{Z}$. Near the vertical boundary,
these equations should be modified as follows:
\[
\eql{4.3} H(-1,s\pm1)\sin\theta_1 - H(0,s) = \mp i\cdot H^{\circ}(-1,s)\cos\theta_1,
\qquad s \notin 2\mathbb{Z}.
\]
Indeed, $X_{[\bvv]}((-\frac12,s\pm\frac12)) = X_{[\bvv]}((0,s)) = H(0,s)$ and
hence (4.3) are nothing but the three-term identities (2.9).

%% =====================================================================
\section*{[Block 2] Enabling construction (Section 4.2) --- printed pp.~2307--2309}
%% =====================================================================

%% --- page 2307 (lower part) ---
\begin{center}
\textbf{4.2. Magnetization $M_m$ in the $(2m)$-th column}
\end{center}

Similarly to Section 3.2, below we rely upon the uniqueness Lemma 4.1 and aim
to construct the values of $X_{[\bvv]}$ on west and east corners (i.e., the
pair of spinors $H,H^{\circ}$) as explicitly as possible. Note that we have
\[
\eql{4.6} H(-2p-1,0) = 0 \ \text{ for } \ p \geqslant m+1,
\qquad H^{\circ}(-2p,0) = 0 \ \text{ for } \ p \leqslant m.
\]

%% --- page 2308 ---
since the spinors defined (on the double cover branching over $\bvv$) by the
symmetry $H_1(-k,-s) := H(-k,s)$, $H_1^{\circ}(-k,-s) := -H^{\circ}(-k,s)$
also satisfy the Cauchy--Riemann equations (4.2), (4.3) and thus must
coincide with $H,H^{\circ}$.

Given $s \geqslant 0$, let $H_s$ denote the semi-infinite vector of the
(real) values $H(-k,s)$, $k \in \mathbb{N}_0$, where we assign zero values
to the indices $s$ such that $s+k \in 2\mathbb{Z}$. Similarly, let
$H_s^{\circ}$ be the vector of the (purely imaginary) values
$H^{\circ}(-k,s)$, $k \in \mathbb{N}$, where we assign zero values to the
indices $s$ such that $s+k \notin 2\mathbb{Z}$. We can write the
harmonicity-type equations (4.4) and (4.5) as
\[
\eql{4.7} H_s = \frac12 C[H_{s-1}+H_{s+1}],\qquad
H_s^{\circ} = \frac12 C^{\circ}[H_{s-1}^{\circ}+H_{s+1}^{\circ}],\qquad s \geqslant 1,
\]
where the self-adjoint operators $C$ and $C^{\circ}$ are given by
\[
C :=
\begin{bmatrix}
0 & \sin\theta_1 & 0 & 0 & \cdots\\
\sin\theta_1 & 0 & \sin\theta_2\cos\theta_1 & 0 & \cdots\\
0 & \sin\theta_2\cos\theta_1 & 0 & \sin\theta_3\cos\theta_2 & \cdots\\
0 & 0 & \sin\theta_3\cos\theta_2 & 0 & \cdots\\
\cdots & \cdots & \cdots & \cdots & \cdots
\end{bmatrix},
\]
\[
C^{\circ} :=
\begin{bmatrix}
0 & \cos\theta_2\sin\theta_1 & 0 & 0 & \cdots\\
\cos\theta_2\sin\theta_1 & 0 & \cos\theta_3\sin\theta_2 & 0 & \cdots\\
0 & \cos\theta_3\sin\theta_2 & 0 & \cos\theta_4\sin\theta_3 & \cdots\\
0 & 0 & \cos\theta_4\sin\theta_3 & 0 & \cdots\\
\cdots & \cdots & \cdots & \cdots & \cdots
\end{bmatrix}.
\]
Let $T(\lambda) := \lambda^{-1}\cdot(1-\sqrt{1-\lambda^2})$. Similarly to
Section 3.2, in order to satisfy the recurrences (4.7) we intend to write
\[
\eql{4.8} H_s := [T(C)]^s H_0, \qquad H_s^{\circ} := [T(C)]^s H_0^{\circ},
\qquad s \geqslant 1.
\]
%% [as printed: (4.8) shows T(C) in BOTH entries.  Copied verbatim.  Reading:
%%  the second T(C) is a misprint for T(C^o) -- Prop. 4.2 and its proof use
%%  [T(C^o)]^s H_0^o, and the sentence on p. 2309 says "the operators T(C) and
%%  T(C^o) in (4.8) are well-defined".]
We now introduce an operator $D$, which plays the key role in the rest of the
analysis:
\[
D := i
\begin{bmatrix}
\cos\theta_1 & 0 & 0 & 0 & \cdots\\
0 & \cos\theta_1\cos\theta_2 & 0 & 0 & \cdots\\
-\sin\theta_1\sin\theta_2 & 0 & \cos\theta_2\cos\theta_3 & 0 & \cdots\\
0 & -\sin\theta_2\sin\theta_3 & 0 & \cos\theta_3\cos\theta_4 & \cdots\\
\cdots & \cdots & \cdots & \cdots & \cdots
\end{bmatrix}.
\]
A straightforward computation gives
\[
\eql{4.9} CD = DC^{\circ},\qquad DD^{*} = I - C^2
\quad\text{and}\quad D^{*}D = I - (C^{\circ})^2.
\]

%% --- page 2309 ---
In particular, this implies that $-I \leqslant C, C^{\circ} \leqslant I$.
Therefore, the operators $T(C)$ and $T(C^{\circ})$ in (4.8) are well-defined
and the vectors $H_s$ and $H_s^{\circ}$ defined by (4.8) are uniformly
bounded as $s \to \infty$. Still, we need to find the vectors $H_0$ and
$H_0^{\circ}$ so that not only the harmonicity-type identities (4.7) for $H$
and $H^{\circ}$ but also the Cauchy--Riemann equations (4.2), (4.3) relating
$H_s$ and $H_s^{\circ}$ are satisfied.

Note that $\operatorname{Ker} D = \{0\}$ while the kernel of $D^{*}$ can be
two-dimensional (more precisely, each of the two operators
$D^{*}_{\mathrm{even}}$ and $D^{*}_{\mathrm{odd}}$ can have a
one-dimensional kernel). Let $D^{*} = U(DD^{*})^{1/2}$ be the polar
decomposition of $D^{*}$, where
\[
\eql{4.10} U := (D^{*}D)^{-1/2}D^{*} = D^{*}(DD^{*})^{-1/2}
\]
is a (partial) isometry. We are now able to formulate the key proposition on
the construction of solutions to (4.2), (4.3) in the upper quadrant.

%% =====================================================================
\section*{[Block 3] COLLECTED THEOREM: Proposition 4.2 and its proof --- printed p.~2309}
%% =====================================================================

\medskip
\noindent\textsc{Proposition 4.2.} --- \emph{Given $H_0 \in \ell^2$, let
$H_0^{\circ} := UH_0$. Then, $H_s := [T(C)]^sH_0$ and
$H_s^{\circ} := [T(C^{\circ})]^sH_0^{\circ}$ are uniformly bounded in
$\ell^2$ and provide a solution to the Cauchy--Riemann equations (4.2), (4.3)
in the upper quadrant.}

\medskip
\noindent\emph{Proof.} --- Since $-I \leqslant C, C^{\circ} \leqslant I$, we
have $0 \leqslant T(C), T(C^{\circ}) \leqslant I$. Therefore, $H_s$ and
$H_s^{\circ}$ are uniformly bounded in $\ell^2$. Moreover, (4.9) and (4.10)
imply that $UC = C^{\circ}U$ and hence
$H_s^{\circ} = [T(C^{\circ})]^sUH_0 = U[T(C)]^sH_0 = UH_s$ for all
$s \geqslant 0$. This allows one to write
\[
\eql{4.11} CH_{s+1}-H_s = -(I-C^2)^{1/2}H_s = -DUH_s = -DH_s^{\circ},
\]
\[
\eql{4.12} H_{s+1}-CH_s = -(I-C^2)^{1/2}H_{s+1} = -DH_{s+1}^{\circ}.
\]
It is not hard to see that these equations are equivalent to the
Cauchy--Riemann identities (4.2), (4.3). Indeed, the first entry of the
vector-valued equation (4.11) or (4.12) (depending on the parity of $s$)
gives the relation (4.3) while the first entry of the other equation gives a
linear combination of (4.3) and (4.2) with $k = 1$. Further, each of the next
entries of (4.11) and (4.12) gives a linear combination of two identities
(4.2) with two consecutive $k$'s. Therefore, for each $s \geqslant 0$ one can
inductively (in $k$) recover all the identities (4.3), (4.2) from (4.11) and
(4.12). \hfill$\square$

\medskip
\noindent\emph{(The paragraph following the proof, still printed p.~2309,
bridges to Theorem 1.1 and is transcribed for completeness.)}

Clearly, the operators $D$ and $U$ can be split into independent components
indexed by odd/even indices, only one of which is relevant for the value of
the magnetization $M_m$ in the even columns $\mathrm{C}_{2m}$, the other
component is responsible for the magnetization in odd columns. In particular,
the relevant block $\Dev$ of the operator $D$ is given by (1.2).

%% --- page 2310 ---
\medskip
\noindent\emph{(Remark 4.3, printed p.~2310, immediately follows and is the
authors' own reading of what Proposition 4.2 constructs.)}

\medskip
\noindent\textsc{Remark 4.3.} --- \emph{In view of the result provided by
Proposition 4.2, the (partial) isometry $\Uev$ can be thought of as a
\textup{discrete Hilbert transform} associated with the Cauchy--Riemann
equations (4.2), (4.3) in the upper quadrant: given the values $H_0$ of the
real part of a ``discrete holomorphic'' function $(H,H^{\circ})$ on the real
line, it returns the boundary values $H_0^{\circ} = \Uev H_0$ of its imaginary
part.}

%% =====================================================================
\section*{[Appendix A] Assembly evidence: the proof of Theorem 1.1 --- printed pp.~2310--2311}
%% =====================================================================

\noindent\emph{(Not part of the collected unit: transcribed to document that
the original target is a two-name reduction, i.e. why it is not the bearer.)}

\medskip
We are now able to prove the main result of this section.

\medskip
\noindent\emph{Proof of Theorem 1.1.} --- Let $H$ and $H^{\circ}$ be the
values of the half-plane observable $X_{[\bvv]}$ on west and east corners,
respectively. Since $H_0$ is a finite vector (see (4.6)), it belongs to
$\ell^2$. Therefore, Lemma 4.1 and Proposition 4.2 imply that
\[
\Dev H_0^{\circ} = \Dev\Uev H_0 = \Sev H_0
= J^{1/2}\,[\,*\ \dots\ *\ M_m\ 0\ 0\ \dots\,]^{\top},
\]
where we use the symbol $*$ to denote unknown entries of the vector $H_0$ and
$M_m$ is its $(m+1)$-th coordinate. On the other hand, note that
\[
-iH^{\circ}(-2m-2,0) = X_{[\bvv]}((-2m-2,0))
= -\mathbb{E}^{+}_{\bHd}\bigl[\sigma_{-2m-\frac52}\bigr] = M_{m+1}.
\]
By definition of the operator $D$ and due to (4.6) one sees that
\[
\Dev H_0^{\circ} = \cos\theta_{2m+1}\cos\theta_{2m+2}
\cdot[\,0\ \dots\ 0\ M_{m+1}\ *\ *\ \dots\,]^{\top}.
\]
Recall that we denote by $P_{m+1}$ the orthogonal projection from $\ell^2$
onto the subspace generated by the first basis vectors $e_1,\dots,e_{m+1}$ of
$\ell^2$. It follows from the considerations given above that
\[
P_{m+1}J^{1/2}P_{m+1} : f_{m+1} = [\,*\ \dots\ *\ 1\,]^{\top}
\mapsto \beta_m\cdot[\,0\ \dots\ 0\ 1\,]^{\top} = \beta_m e_{m+1}
\]
for a certain vector $f_{m+1} = P_{m+1}f_{m+1}$ such that
$\langle f_{m+1}, e_{m+1}\rangle = 1$, where
\[
\beta_m := \cos\theta_{2m+1}\cos\theta_{2m+2}\cdot M_{m+1}/M_m.
\]
In particular, if we denote by $e_1',e_2',\dots$ the orthogonalization of
the vectors $e_1,e_2,\dots$ with respect to the scalar product
$\langle\,\cdot\,,J^{1/2}\,\cdot\,\rangle$, then
\[
\langle e_{m+1}',J^{1/2}e_{m+1}'\rangle
= \langle e_{m+1},J^{1/2}f_{m+1}\rangle = \beta_m
\]
and hence
\begin{align*}
\det P_{m+1}J^{1/2}P_{m+1}
&= \det[\langle e_p',J^{1/2}e_q'\rangle]_{p,q=1}^{m+1}\\
&= \prod_{k=0}^{m}\beta_k = M_{m+1}\cdot\prod_{k=1}^{2m+2}\cos\theta_k,
\end{align*}
where we also used the fact that $M_0 = 1$; note that this computations does
not require any modification in the case $m = 0$ (when dealing with

%% --- page 2311 ---
the magnetization in even columns). This gives the second formula for $M_m$
in (1.5).

To prove that $M_m$ also equals $|\det P_m\Uev P_m|$, note that
\[
(\Dev D_{\mathrm{even}}^{*})^{1/2} = \Dev\Uev
\quad\text{and}\quad P_m\Dev = P_m\Dev P_m,
\]
which implies
\begin{align*}
\det P_m(\Dev D_{\mathrm{even}}^{*})^{1/2}P_m
&= |\det P_m\Uev P_m|\cdot|\det P_m\Dev P_m|\\
&= |\det P_m\Uev P_m|\cdot\prod_{k=1}^{2m}\cos\theta_k.
\end{align*}
Finally, to prove the last identity in (1.5), note that
\[
\det P_mJ^{1/2}P_m
= \frac{\det[\langle J^{1/2}f_p,f_q\rangle]_{p,q=1}^{m}}
{\det[\langle f_p,f_q\rangle]_{p,q=1}^{m}}
\]
for all bases $f_1,\dots,f_m$ of the $m$-dimensional space
$\operatorname{Ran}P_m$. Choosing the basis $1,\lambda,\dots,\lambda^{m-1}$ in
the spectral representation of the operator $J$ in $L^2(\nu_J(d\lambda))$ one
obtains the identity
\[
\det P_mJ^{1/2}P_m = \frac{H_m[\lambda^{1/2}\nu_J]}{H_m[\nu_J]}
\quad\text{and, similarly,}\quad
\det P_mJP_m = \frac{H_m[\lambda\nu_J]}{H_m[\nu_J]}
\]
As $\det P_mJP_m = [\det P_m\Dev P_m]^2
= \prod_{k=1}^{2n}\cos^2\theta_k$, this completes the proof.
\hfill$\square$
%% [as printed: the last product is printed with "2n" (not "2m"); copied verbatim.]

%% =====================================================================
\section*{[Appendix B] Co-cited companion: Lemma 4.1 and its proof --- printed pp.~2305--2307}
%% =====================================================================

\noindent\emph{(This is the second name in the deferral ``Lemma 4.1 and
Proposition 4.2 imply that''. It is transcribed in full so that the bearer of
Theorem 1.1 can be re-judged.)}

\medskip
%% --- page 2305 ---
\noindent\textsc{Lemma 4.1.} --- \emph{The spinors $H,H^{\circ}$ defined in
$\bHd$ and branching over $\bvv$ are uniquely determined by the following
conditions: uniform boundedness, Cauchy--Riemann identities (4.2), boundary
relations (4.3), and the value (4.1) of $H$ near $\bvv$.}

\medskip
\noindent\emph{Proof.} --- Taking the difference of two solutions, assume
that $H,H^{\circ}$ are uniformly bounded, satisfy (4.2), (4.3) and that
$H(-2m-1,0) = 0$. Recall that Proposition 2.4 gives the harmonicity-type
identity
\begin{align*}
\eql{4.4} H(-k,s) &= \frac12\sin\theta_{k+1}\cos\theta_k
\cdot[H(-k-1,s+1)+H(-k-1,s-1)]\\
&\quad + \frac12\sin\theta_k\cos\theta_{k-1}
\cdot[H(-k+1,s+1)+H(-k+1,s-1)]
\end{align*}
at all west corners $c = (-k+\frac12,s)$ with $k \geqslant 2$ except in the
case $k = -2m-1$, $s = 0$ (i.e., at the west corner located near the branching
$\bvv$). Moreover, due to the boundary relations (4.3), exactly the same
identity holds for $k = 0,1$ (recall that we formally set $\theta_0 := 0$). In
its turn, the function $H^{\circ}$ satisfies the identities
\begin{align*}
\eql{4.5} H^{\circ}(-k,s) &= \frac12\cos\theta_{k+1}\sin\theta_k
\cdot[H^{\circ}(-k-1,s+1)+H^{\circ}(-k-1,s-1)]\\
&\quad + \frac12\cos\theta_k\sin\theta_{k-1}
\cdot[H^{\circ}(-k+1,s+1)+H^{\circ}(-k+1,s-1)]
\end{align*}
at all east corners $d = (-k+\frac12,s)$, including the one located near the
branching $\bvv$ (in the latter case the proof of Proposition 2.4 works
verbatim due to the fact that $H(-2m-1,0) = 0$). Both (4.4) and (4.5) can be
rewritten as true discrete harmonicity properties if one passes from $H$

%% --- page 2306 ---
and $H^{\circ}$ to the functions
\[
\Htil(-k,s) := \varrho_k\cdot H(-k,s),\qquad
\Htilo := \varrho^{\circ}_k\cdot H^{\circ}(-k,s),
\]
%% [as printed: the argument "(-k, s)" is missing on the left-hand side of the
%%  second definition; copied verbatim.]
\[
\varrho_k := \prod_{j=1}^{k}(\sin\theta_j/\cos\theta_{j-1}),\qquad
\varrho^{\circ}_k := \prod_{j=2}^{k}(\cos\theta_j/\sin\theta_{j-1}),
\]
recall that we set $\Htilo(0,s) = H^{\circ}(0,s) := 0$ on the vertical axes.
Let $Z_n = (K_n,S_n)$ (respectively, $Z_n^{\circ} = (K_n^{\circ},S_n)$) be the
nearest-neighbor random walk on west (respectively, east) corners, with jump
probabilities $(\frac12,\frac12)$ for the process $S_n$ and
$(\cos^2\theta_k,\sin^2\theta_k)$ for the process $K_n$ (respectively,
$(\sin^2\theta_k,\cos^2\theta_k)$ for the process $K_n^{\circ}$); see Figure
4.2. Note that the walk $Z_n$ on west corners is reflected from the vertical
axes while the walk $Z_n^{\circ}$ on east corners is absorbed there.
%% [FIGURE 4.2 not transcribed: two horizontal chains of corners C_{k+1}, C_k,
%%  C_{k-1}, ... -> ... -> C_2, C_1, C_0, the upper chain carrying the up/down
%%  transition probabilities sin^2 theta_j / cos^2 theta_j, the lower chain the
%%  swapped ones, the lower-right extremity ending in a self-loop at C_0.]
\begin{quote}
\small Figure 4.2. For appropriately chosen prefactors $\varrho_k$ and
$\varrho^{\circ}_k$, the identities (4.4), (4.5) (coming from Proposition 2.4)
can be written as the discrete harmonicity property of functions
$\varrho_kH(-k,s)$ and $\varrho^{\circ}_kH^{\circ}(-k,s)$ with respect to
random walks having the indicated transition probabilities in the horizontal
direction (and $\frac12$ in the vertical one). The first random walk (on
$\triangleright$) is reflected from the imaginary axis while the second (on
$\triangleleft$) is absorbed there.
\end{quote}
It follows from (4.4) that the stochastic process $\Htil(Z_n)$ is a
martingale, when equipped with the canonical filtration, until the first time
when $Z_n$ hits the west corner $(-2m-1,0)$ located near the branching, recall
that $\Htil(-2m-1,0) = 0$. Similarly, (4.5) implies that the process
$\Htilo(Z_n^{\circ})$ is a martingale until the first time when $Z_n^{\circ}$
hits the imaginary axis, recall that $\Htilo = 0$ there. As we show below,
depending on the behavior of $\varrho_k$ and $\varrho^{\circ}_k$ as
$k \to \infty$, the optional stopping theorem allows to conclude that either
$\Htil$ or $\Htilo$ vanishes identically. Once the identity
$\Htil \equiv 0$ (resp., $\Htilo \equiv 0$) is proven, the equations (4.2),
(4.3) and the fact that $H^{\circ}$ vanishes

%% --- page 2307 ---
on the imaginary axis (resp., $\Htil$ vanishes at the point $(-2m-1,0)$)
imply that $\Htilo \equiv 0$ (resp., $\Htil \equiv 0$) too.

Recall that the functions $H$ and $H^{\circ}$ are uniformly bounded and note
that $\varrho_k\varrho^{\circ}_k = (\cos\theta_1)^{-1}\sin\theta_k\cos\theta_k
= O(1)$ as $k \to \infty$. It follows from the maximum principle that
\begin{itemize}
\item the function $\Htil$ is uniformly bounded unless $\varrho_k \to \infty$
as $k \to \infty$;
\item the function $\Htilo$ is uniformly bounded unless
$\varrho^{\circ}_k \to \infty$ as $k \to \infty$.
\end{itemize}
We have three cases to consider separately.
\begin{itemize}
\item Let $\liminf_{k\to\infty}\varrho_k = 0$, in particular this implies that
$\Htil$ is uniformly bounded. The optional stopping theorem applied to the
martingale $\Htil(Z_n)$ and the fact that a one-dimensional random walk on
$-\mathbb{N}_0$ reflected at $0$ almost surely takes arbitrary large
(negative) values imply that $H \equiv 0$.
\item Let $\liminf_{k\to\infty}\varrho^{\circ}_k = 0$. A similar argument
applied to the martingale $\Htilo(Z_n^{\circ})$ (recall that $\Htilo$ vanishes
on the imaginary axis) shows that $\Htilo \equiv 0$.
\item Let both sequences $\varrho_k$ and $\varrho^{\circ}_k$ be uniformly
bounded from below as $k \to \infty$. Since
$\varrho_k\varrho^{\circ}_k = (\cos\theta_1)^{-1}\sin\theta_k\cos\theta_k$,
these sequences are also uniformly bounded from above and the parameters
$\theta_k$, $k \geqslant 1$, stay away from $0$. In this case it is easy to
see that the process $K_n^{\circ}$ hits $0$ almost surely (i.e., that the
random walk $Z_n^{\circ}$ hits the imaginary axis almost surely). Indeed, the
probability $p_k^{\circ}$ to hit $0$ starting from $-k$ satisfies the
recurrence
\[
p_k^{\circ}-p_{k+1}^{\circ} = \cot^2\theta_k\cdot(p_{k-1}^{\circ}-p_k^{\circ})
= \dots = \varrho_{k+1}^{-2}\sin^2\theta_{k+1}\cdot(1-p_1^{\circ}),
\]
which is only possible if $p_1^{\circ} = 1$ since the factors
$\varrho_{k+1}/\sin\theta_{k+1}$ are uniformly bounded. We conclude as before
by applying the optional stopping theorem to the uniformly bounded martingale
$\Htil(Z_n^{\circ})$.
%% [as printed: this last martingale is printed H~(...) without the ring, while
%%  the case being closed is the one about H~^o; copied verbatim.]
\end{itemize}
The proof is complete. \hfill$\square$

%% =====================================================================
\section*{[Appendix C] Flagged source oddities (all copied verbatim above)}
%% =====================================================================
\begin{itemize}
\item (4.2) printed as ``$= ,\pm i\cdot[\dots]$'' (stray comma after the equality sign).
\item (4.8) prints $[T(C)]^sH_0^{\circ}$ where $[T(C^{\circ})]^sH_0^{\circ}$ is meant
      (Prop. 4.2; the sentence ``the operators $T(C)$ and $T(C^{\circ})$ in (4.8)'').
\item p. 2306: the definition of $\Htilo$ omits its argument $(-k,s)$.
\item p. 2307, last bullet: the martingale is printed $\Htil(Z_n^{\circ})$, but the
      identity being closed there is the one about $\Htilo$.
\item p. 2310: ``note that this computations does not require'' (as printed).
\item p. 2311: the final product is printed $\prod_{k=1}^{2n}\cos^2\theta_k$ (``$2n$'' for ``$2m$'').
\end{itemize}

\end{document}
